For a chord $UV$ in a circle with centre $O$ and radius $R$, the perpendicular from $O$ bisects the chord at $M$. Pythagoras gives $O M^2 = R^2 - (\dfrac{U V}{2})^2$. Thus equal chord lengths give equal distances from the centre. Applying this to $AC=BD=CE$ proves the three chord lines are equidistant from $O$. Their common perpendicular distance is $\boxed{\sqrt{R^2 - \dfrac{\ell^2}{4}}}$ if their common length is $\ell$.
\begin{asy}
import graph;
usepackage("amsmath");
import geometry;
size(7cm);
pair O=(0,0);for(int i=0;i<5;++i){pair P=2dir(90+72i);label("$"+substr("ABCDE",i,1)+"$",1.15P);}draw(circle(O,2));draw(2dir(90)--2dir(234));draw(2dir(162)--2dir(306));draw(2dir(234)--2dir(378));dot(O);label("$O$",O,(0,-1));
\end{asy}

